\documentclass{amsart}
% For dealing with references we use the comment environment
\usepackage{verbatim}
\newenvironment{reference}{\comment}{\endcomment}
%\newenvironment{reference}{}{}
\newenvironment{slogan}{\comment}{\endcomment}
\newenvironment{history}{\comment}{\endcomment}

% For commutative diagrams we use Xy-pic
\usepackage[all]{xy}

% We use 2cell for 2-commutative diagrams.
\xyoption{2cell}
\UseAllTwocells

% We use multicol for the list of chapters between chapters
\usepackage{multicol}

% This is generally recommended for better output
\usepackage{lmodern}
\usepackage[T1]{fontenc}

% For cross-file-references

% Package for hypertext links:
\usepackage{hyperref}

% For any local file, say "hello.tex" you want to link to please

% Theorem environments.
%
\theoremstyle{plain}
\newtheorem{theorem}[subsection]{Theorem}
\newtheorem{proposition}[subsection]{Proposition}
\newtheorem{lemma}[subsection]{Lemma}

\theoremstyle{definition}
\newtheorem{definition}[subsection]{Definition}
\newtheorem{example}[subsection]{Example}
\newtheorem{exercise}[subsection]{Exercise}
\newtheorem{situation}[subsection]{Situation}

\theoremstyle{remark}
\newtheorem{remark}[subsection]{Remark}
\newtheorem{remarks}[subsection]{Remarks}

\numberwithin{equation}{subsection}

% Macros
%
\def\lim{\mathop{\mathrm{lim}}\nolimits}
\def\colim{\mathop{\mathrm{colim}}\nolimits}
\def\Spec{\mathop{\mathrm{Spec}}}
\def\Hom{\mathop{\mathrm{Hom}}\nolimits}
\def\Ext{\mathop{\mathrm{Ext}}\nolimits}
\def\SheafHom{\mathop{\mathcal{H}\!\mathit{om}}\nolimits}
\def\SheafExt{\mathop{\mathcal{E}\!\mathit{xt}}\nolimits}
\def\Sch{\mathit{Sch}}
\def\Mor{\mathop{\mathrm{Mor}}\nolimits}
\def\Ob{\mathop{\mathrm{Ob}}\nolimits}
\def\Sh{\mathop{\mathit{Sh}}\nolimits}
\def\NL{\mathop{N\!L}\nolimits}
\def\CH{\mathop{\mathrm{CH}}\nolimits}
\def\proetale{{pro\text{-}\acute{e}tale}}
\def\etale{{\acute{e}tale}}
\def\QCoh{\mathit{QCoh}}
\def\Ker{\mathop{\mathrm{Ker}}}
\def\Im{\mathop{\mathrm{Im}}}
\def\Coker{\mathop{\mathrm{Coker}}}
\def\Coim{\mathop{\mathrm{Coim}}}

% Boxtimes
%
\DeclareMathSymbol{\boxtimes}{\mathbin}{AMSa}{"02}

%
% Macros for moduli stacks/spaces
%
\def\QCohstack{\mathcal{QC}\!\mathit{oh}}
\def\Cohstack{\mathcal{C}\!\mathit{oh}}
\def\Spacesstack{\mathcal{S}\!\mathit{paces}}
\def\Quotfunctor{\mathrm{Quot}}
\def\Hilbfunctor{\mathrm{Hilb}}
\def\Curvesstack{\mathcal{C}\!\mathit{urves}}
\def\Polarizedstack{\mathcal{P}\!\mathit{olarized}}
\def\Complexesstack{\mathcal{C}\!\mathit{omplexes}}
% \Pic is the operator that assigns to X its picard group, usage \Pic(X)
% \Picardstack_{X/B} denotes the Picard stack of X over B
% \Picardfunctor_{X/B} denotes the Picard functor of X over B
\def\Pic{\mathop{\mathrm{Pic}}\nolimits}
\def\Picardstack{\mathcal{P}\!\mathit{ic}}
\def\Picardfunctor{\mathrm{Pic}}
\def\Deformationcategory{\mathcal{D}\!\mathit{ef}}

\setlength{\emergencystretch}{3em}
\begin{document}
\title[Stacks Project, Tag 0FBD]{454 --- Localized Chern classes satisfy triangle additivity under rank and representability hypotheses}
\maketitle
\noindent Copyright (C) 2005--2025 Johan de Jong. Stacks Project authors. Source: \href{https://stacks.math.columbia.edu/tag/0FBD}{Tag 0FBD}.

\noindent Editorial difficulty note (not author text). Difficulty H3: The proof must establish multiplicativity in the supported bivariant group, including all listed representability alternatives. It deforms a triangle map by a scalar over the affine line to compare the cone with the direct sum, then combines two graph constructions over a common fibre product to prove the product formula there..

\noindent Editorial provenance note (not author text). History: excerpt from revision \texttt{a0444\allowbreak 6e57e\allowbreak c1fbc\allowbreak 252a8\allowbreak 71afc\allowbreak ec775\allowbreak 2fb28\allowbreak 07b14}, chow.tex, lines 11329--11479. Reference presentation was adapted; original-author context and core support blocks were retained or added. The mathematical wording is unchanged. Licensed under GNU FDL 1.2 or later; no invariant sections or cover texts. See ../licenses/COPYING. Context blocks reproduce author definitions and supporting results; they are not additional counted problems.

\begin{definition}
\label{definition-bivariant-class}
\begin{reference}
Similar to \href{https://stacks.math.columbia.edu/bibliography}{Bibliography: F (Definition 17.1)}
\end{reference}
Let $(S, \delta)$ be as in Situation \ref{situation-setup}.
Let $f : X \to Y$ be a morphism of schemes locally of finite type over $S$.
Let $p \in \mathbf{Z}$.
A {\it bivariant class $c$ of degree $p$ for $f$} is given by a rule
which assigns to every locally of finite type morphism $Y' \to Y$
and every $k$ a map
$$
c \cap - : \CH_k(Y') \longrightarrow \CH_{k - p}(X')
$$
where $X' = Y' \times_Y X$, satisfying the following conditions
\begin{enumerate}
\item if $Y'' \to Y'$ is a proper, then
$c \cap (Y'' \to Y')_*\alpha'' = (X'' \to X')_*(c \cap \alpha'')$
for all $\alpha''$ on $Y''$ where $X'' = Y'' \times_Y X$,
\item if $Y'' \to Y'$ is flat locally of finite type of
fixed relative dimension, then
$c \cap (Y'' \to Y')^*\alpha' = (X'' \to X')^*(c \cap \alpha')$
for all $\alpha'$ on $Y'$, and
\item if $(\mathcal{L}', s', i' : D' \to Y')$ is as in
Definition \ref{definition-gysin-homomorphism}
with pullback $(\mathcal{N}', t', j' : E' \to X')$ to $X'$,
then we have $c \cap (i')^*\alpha' = (j')^*(c \cap \alpha')$
for all $\alpha'$ on $Y'$.
\end{enumerate}
The collection of all bivariant classes of degree $p$ for $f$ is
denoted $A^p(X \to Y)$.
\end{definition}

\begin{situation}
\label{situation-setup}
Here $S$ is a locally Noetherian, and universally catenary scheme.
Moreover, we assume $S$ is endowed with a dimension function
$\delta : S \longrightarrow \mathbf{Z}$.
\end{situation}

\begin{definition}
\label{definition-gysin-homomorphism}
Let $(S, \delta)$ be as in Situation \ref{situation-setup}.
Let $X$ be locally of finite type over $S$.
Let $(\mathcal{L}, s)$ be a pair consisting of an invertible
sheaf and a global section $s \in \Gamma(X, \mathcal{L})$.
Let $D = Z(s)$ be the zero scheme of $s$, and
denote $i : D \to X$ the closed immersion.
We define, for every integer $k$, a {\it Gysin homomorphism}
$$
i^* : Z_{k + 1}(X) \to \CH_k(D).
$$
by the following rules:
\begin{enumerate}
\item Given an integral closed subscheme $W \subset X$ with
$\dim_\delta(W) = k + 1$ we define
\begin{enumerate}
\item if $W \not \subset D$, then $i^*[W] = [D \cap W]_k$ as a
$k$-cycle on $D$, and
\item if $W \subset D$, then
$i^*[W] = i'_*(c_1(\mathcal{L}|_W) \cap [W])$,
where $i' : W \to D$ is the induced closed immersion.
\end{enumerate}
\item For a general $(k + 1)$-cycle $\alpha = \sum n_j[W_j]$
we set
$$
i^*\alpha = \sum n_j i^*[W_j]
$$
\item If $D$ is an effective Cartier divisor, then we denote
$D \cdot \alpha = i_*i^*\alpha$ the pushforward of the class $i^*\alpha$
to a class on $X$.
\end{enumerate}
\end{definition}

\begin{situation}
\label{situation-loc-chern}
Let $(S, \delta)$ be as in Situation \ref{situation-setup}. Let $X$ be a scheme
locally of finite type over $S$. Let $i : Z \to X$ be a closed immersion.
Let $E \in D(\mathcal{O}_X)$ be an object. Let $p \geq 0$. Assume
\begin{enumerate}
\item $E$ is a perfect object of $D(\mathcal{O}_X)$,
\item the restriction $E|_{X \setminus Z}$ is zero, resp.\ isomorphic to a
finite locally free $\mathcal{O}_{X \setminus Z}$-module of rank $< p$
sitting in cohomological degree $0$, and
\item at least one\footnote{Please ignore this technical condition on a
first reading; see discussion in Remark \href{https://stacks.math.columbia.edu/tag/0GUK}{remark loc chern (Tag 0GUK)}.}
of the following is true:
(a) $X$ is quasi-compact,
(b) $X$ has quasi-compact irreducible components,
(c) there exists a locally bounded complex of finite locally free
$\mathcal{O}_X$-modules representing $E$, or
(d) there exists a morphism $X \to X'$ of schemes locally of finite type
over $S$ such that $E$ is the pullback of a perfect object on $X'$ and
the irreducible components of $X'$ are quasi-compact.
\end{enumerate}
\end{situation}

\begin{remark}
\label{remark-loc-chern-classes}
In Situation \ref{situation-loc-chern} it is convenient to define
$$
c^{(p)}(Z \to X, E) = 1 + c_1(E) + \ldots + c_{p - 1}(E) +
c_p(Z \to X, E) + c_{p + 1}(Z \to X, E) + \ldots
$$
as an element of the algebra $A^{(p)}(Z \to X)$ considered in
Remark \href{https://stacks.math.columbia.edu/tag/0FA0}{remark ring loc classes (Tag 0FA0)}.
\end{remark}

\begin{lemma}
\label{lemma-additivity-loc-chern-c}
Let $(S, \delta)$ be as in Situation \ref{situation-setup}.
Let $X$ be locally of finite type over $S$. Let $Z \to X$ be
a closed immersion. Let
$$
E_1 \to E_2 \to E_3 \to E_1[1]
$$
be a distinguished triangle of perfect objects in $D(\mathcal{O}_X)$.
Assume
\begin{enumerate}
\item the restrictions $E_1|_{X \setminus Z}$ and $E_3|_{X \setminus Z}$
are isomorphic to finite locally free $\mathcal{O}_{X \setminus Z}$-modules
of rank $< p_1$ and $< p_3$ placed in degree $0$, and
\item at least one of the following is true:
(a) $X$ is quasi-compact,
(b) $X$ has quasi-compact irreducible components,
(c) $E_3 \to E_1[1]$ can be represented by a map of locally
bounded complexes of finite locally free $\mathcal{O}_X$-modules, or
(d) there exists an envelope $f : Y \to X$ such that $Lf^*E_3 \to Lf^*E_1[1]$
can be represented by a map of locally bounded complexes of
finite locally free $\mathcal{O}_Y$-modules.
\end{enumerate}
With notation as in Remark \ref{remark-loc-chern-classes} we have
$$
c^{(p_1 + p_3)}(Z \to X, E_2) = c^{(p_1)}(Z \to X, E_1)c^{(p_3)}(Z \to X, E_3)
$$
in $A^{(p_1 + p_3)}(Z \to X)$.
\end{lemma}

\begin{proof}
Observe that the assumptions imply that $E_2|_{X \setminus Z}$ is zero,
resp.\ isomorphic to a finite locally free $\mathcal{O}_{X \setminus Z}$-module
of rank $< p_1 + p_3$. Thus the statement makes sense.

\medskip\noindent
Let $f : Y \to X$ be an envelope. Expanding the left and right hand sides
of the formula in the statement of the lemma we see that we have to prove
some equalities of classes in $A^*(X)$ and in $A^*(Z \to X)$. By the
uniqueness in Lemma \href{https://stacks.math.columbia.edu/tag/0GUC}{lemma envelope bivariant (Tag 0GUC)} it suffices to prove the
corresponding relations in $A^*(Y)$ and $A^*(Z \to Y)$. Since moreover
the construction of the classes involved is compatible with base change
(Lemma \href{https://stacks.math.columbia.edu/tag/0FB3}{lemma base change loc chern (Tag 0FB3)}) we may replace $X$ by $Y$
and the distinguished triangle by its pullback.

\medskip\noindent
In the proof of Lemma \href{https://stacks.math.columbia.edu/tag/0F9F}{lemma additivity on perfect (Tag 0F9F)} we have
seen that conditions (2)(a), (2)(b), and (2)(c) imply condition
(2)(d). Combined with the discussion in the previous paragraph we
reduce to the case discussed in the next paragraph.

\medskip\noindent
Let $\varphi^\bullet : \mathcal{E}_3^\bullet[-1] \to \mathcal{E}_1^\bullet$
be a map of locally bounded complexes of finite locally free
$\mathcal{O}_X$-modules representing the map $E_3[-1] \to E_1$
in the derived category. Consider the scheme
$X' = \mathbf{A}^1 \times X$ with projection
$g : X' \to X$. Let $Z' = g^{-1}(Z) = \mathbf{A}^1 \times Z$.
Denote $t$ the coordinate on $\mathbf{A}^1$. Consider the cone
$\mathcal{C}^\bullet$ of the map of complexes
$$
t g^*\varphi^\bullet :
g^*\mathcal{E}_3^\bullet[-1]
\longrightarrow
g^*\mathcal{E}_1^\bullet
$$
over $X'$. We obtain a distinguished triangle
$$
g^*\mathcal{E}_1^\bullet \to \mathcal{C}^\bullet \to
g^*\mathcal{E}_3^\bullet \to g^*\mathcal{E}_1^\bullet[1]
$$
where the first three terms form a termwise split short exact
sequence of complexes. Clearly $\mathcal{C}^\bullet$ is a
bounded complex of finite locally free $\mathcal{O}_{X'}$-modules
whose restriction to $X' \setminus Z'$ is isomorphic to a
finite locally free
$\mathcal{O}_{X' \setminus Z'}$-module of rank $< p_1 + p_3$
placed in degree $0$. Thus we have the localized Chern classes
$$
c_p(Z' \to X', \mathcal{C}^\bullet) \in A^p(Z' \to X')
$$
for $p \geq p_1 + p_3$. For any $\alpha \in \CH_k(X)$ consider
$$
c_p(Z' \to X', \mathcal{C}^\bullet) \cap g^*\alpha
\in \CH_{k + 1 - p}(\mathbf{A}^1 \times X)
$$
If we restrict to $t = 0$, then the map $t g^*\varphi^\bullet$
restricts to zero and $\mathcal{C}^\bullet|_{t = 0}$
is the direct sum of $\mathcal{E}_1^\bullet$ and $\mathcal{E}_3^\bullet$.
By compatibility of localized Chern classes with base change
(Lemma \href{https://stacks.math.columbia.edu/tag/0FB3}{lemma base change loc chern (Tag 0FB3)}) we conclude that
$$
i_0^* \circ c^{(p_1 + p_3)}(Z' \to X', \mathcal{C}^\bullet) \circ g^* =
c^{(p_1 + p_2)}(Z \to X, E_1 \oplus E_3)
$$
in $A^{(p_1 + p_3)}(Z \to X)$. On the other hand, if we restrict to $t = 1$,
then the map $t g^*\varphi^\bullet$
restricts to $\varphi$ and $\mathcal{C}^\bullet|_{t = 1}$
is a bounded complex of finite locally free modules representing $E_2$.
We conclude that
$$
i_1^* \circ c^{(p_1 + p_3)}(Z' \to X', \mathcal{C}^\bullet) \circ g^* =
c^{(p_1 + p_2)}(Z \to X, E_2)
$$
in $A^{(p_1 + p_3)}(Z \to X)$. Since $i_0^* = i_1^*$ by definition of
rational equivalence (more precisely this follows from the formulae in
Lemma \href{https://stacks.math.columbia.edu/tag/0F96}{lemma linebundle formulae (Tag 0F96)}) we conclude that
$$
c^{(p_1 + p_2)}(Z \to X, E_2) = c^{(p_1 + p_2)}(Z \to X, E_1 \oplus E_3)
$$
This reduces us to the case discussed in the next paragraph.

\medskip\noindent
Assume $E_2 = E_1 \oplus E_3$ and the triples $(X, Z, E_i)$
are as in Situation \ref{situation-loc-chern}.
For $i = 1, 3$ let
$$
b_i : W_i \to \mathbf{P}^1_X
\quad\text{and}\quad
Q_i
\quad\text{and}\quad
T'_i \subset T_i \subset W_{i, \infty}
$$
be as in the proof of Lemma \href{https://stacks.math.columbia.edu/tag/0FB2}{lemma independent loc chern (Tag 0FB2)}.
By definition
$$
c_p(Z \to X, E_i) = c'_p(Q_i)
$$
where the right hand side is the bivariant class constructed in
Lemma \href{https://stacks.math.columbia.edu/tag/0F9K}{lemma localized chern pre (Tag 0F9K)} using $W_i, b_i, Q_i, T'_i$.
Set $W = W_1 \times_{b_1, \mathbf{P}^1_X, b_2} W_2$ and consider
the cartesian diagram
$$
\xymatrix{
W \ar[d]_{g_1} \ar[rd]^b \ar[r]_{g_3} & W_3 \ar[d]^{b_3} \\
W_1 \ar[r]^{b_1} & \mathbf{P}^1_X
}
$$
Of course $b^{-1}(\mathbf{A}^1)$ maps isomorphically to $\mathbf{A}^1_X$.
Observe that $T' = g_1^{-1}(T'_1) \cap g_2^{-1}(T'_2)$ still contains
all the points of $W_\infty$ lying over $X \setminus Z$.
By Lemma \href{https://stacks.math.columbia.edu/tag/0FAU}{lemma localized chern pre independent (Tag 0FAU)} we may use
$W$, $b$, $g_i^*\mathcal{Q}_i$, and
$T'$ to construct $c_p(Z \to X, E_i)$ for $i = 1, 3$.
Also, by the stronger independence given in
Lemma \href{https://stacks.math.columbia.edu/tag/0FE8}{lemma independent loc chern bQ (Tag 0FE8)} we may use
$W$, $b$, $g_1^*Q_1 \oplus g_3^*Q_3$, and $T'$
to compute the classes $c_p(Z \to X, E_2)$.
Thus the desired equality follows from
Lemma \href{https://stacks.math.columbia.edu/tag/0FAY}{lemma localized chern pre sum c (Tag 0FAY)}.
\end{proof}
\end{document}
